%% ma: L12-B16-0019
%% lop: 12
%% chuong: 5
%% bai: 16
%% dang: 12.16.8
%% loai: TN4
%% muc: TH
%% dap_an: "A"
%% nguon: "Ảnh giáo viên gửi 04/11/2026 (Huỳnh Văn Ánh, Chuyên đề V - Phương pháp tọa độ trong không gian), Đề test theo bài, Phần I Câu 10"
%% kien_thuc: "Bài 16 (góc giữa hai đường thẳng), Bài 15 (vectơ chỉ phương từ phương trình chính tắc và tham số)"
%% trang_thai: chinh-thuc
%% ghi_chu: "Đáp án gốc A khớp (sympy). Đã duyệt 04/11/2026."
\begin{ex}
Trong không gian với hệ tọa độ $Oxyz$, tính góc giữa hai đường thẳng
\[d_1\colon\dfrac{x-1}{1}=\dfrac{y+5}{-1}=\dfrac z4\quad\text{và}\quad d_2\colon\begin{cases}x=3+2t\\y=-3\\z=2-2t.\end{cases}\]
\choice{\True $60^\circ$}{$90^\circ$}{$30^\circ$}{$45^\circ$}
\loigiai{$d_1$ có vectơ chỉ phương $\vec u_1=(1;-1;4)$, $d_2$ có vectơ chỉ phương $\vec u_2=(2;0;-2)$.

Suy ra $\cos(d_1,d_2)=\dfrac{|\vec u_1\cdot\vec u_2|}{|\vec u_1|\cdot|\vec u_2|}=\dfrac{|2+0-8|}{3\sqrt2\cdot2\sqrt2}=\dfrac6{12}=\dfrac12$. Vậy $(d_1,d_2)=60^\circ$. Chọn A.}
\end{ex}
